% F081: precise mathematical teaching diagram.
\begin{tikzpicture}[x=1cm,y=1cm,>=stealth,font=\small]
\draw[rounded corners=3pt,gray!40] (0,0) rectangle (14.2,8.5);
\begin{scope}[yshift=1cm]
\node at (7.1,6.9) {两端等高：上拱、平走、下凹};
\node[] at (2.6,6.3) {先上升，再下降};
\begin{scope}[shift={(2.6,3.9)},x=1.35cm,y=1.4cm]
\draw[->,gray!65] (-1.35,0)--(1.35,0) node[right] {$x$};
\draw[->,gray!65] (0,-1.45)--(0,1.2) node[left] {$y$};
\draw[blue!75!black,line width=1.15pt,,domain=-1:1,samples=160] plot (\x,{1-(\x)*(\x)});
\fill[blue!75!black] (-1,0) circle (1.8pt);\node[above left] at (-1,0) {$A$};
\fill[blue!75!black] (1,0) circle (1.8pt);\node[above right] at (1,0) {$B$};
\draw[orange!85!black,line width=1.2pt] (-0.3,1)--(0.3,1);
\fill[orange!85!black] (0,1) circle (1.8pt);
\end{scope}
\node[] at (2.6,1.25) {内部最高处：$f^{\prime}(c)=0$};
\node[] at (7.1,6.3) {全程保持等高};
\begin{scope}[shift={(7.1,3.9)},x=1.35cm,y=1.4cm]
\draw[->,gray!65] (-1.35,0)--(1.35,0) node[right] {$x$};
\draw[->,gray!65] (0,-1.45)--(0,1.2) node[left] {$y$};
\draw[blue!75!black,line width=1.15pt,,domain=-1:1,samples=160] plot (\x,{0});
\fill[blue!75!black] (-1,0) circle (1.8pt);\node[above left] at (-1,0) {$A$};
\fill[blue!75!black] (1,0) circle (1.8pt);\node[above right] at (1,0) {$B$};
\draw[orange!85!black,line width=1.2pt] (-0.3,0)--(0.3,0);
\fill[orange!85!black] (0,0) circle (1.8pt);
\end{scope}
\node[] at (7.1,1.25) {内部每一点：$f^{\prime}(c)=0$};
\node[] at (11.6,6.3) {先下降，再上升};
\begin{scope}[shift={(11.6,3.9)},x=1.35cm,y=1.4cm]
\draw[->,gray!65] (-1.35,0)--(1.35,0) node[right] {$x$};
\draw[->,gray!65] (0,-1.45)--(0,1.2) node[left] {$y$};
\draw[blue!75!black,line width=1.15pt,,domain=-1:1,samples=160] plot (\x,{(\x)*(\x)-1});
\fill[blue!75!black] (-1,0) circle (1.8pt);\node[above left] at (-1,0) {$A$};
\fill[blue!75!black] (1,0) circle (1.8pt);\node[above right] at (1,0) {$B$};
\draw[orange!85!black,line width=1.2pt] (-0.3,-1)--(0.3,-1);
\fill[orange!85!black] (0,-1) circle (1.8pt);
\end{scope}
\node[] at (11.6,1.25) {内部最低处：$f^{\prime}(c)=0$};
\end{scope}
\node[text width=13cm,align=center] at (7.1,.9) {三种典型情形均有水平切线；前提仍是闭区间连续、开区间可导。};
\end{tikzpicture}
